\section{All mixed finite amplitude covariances with the original two global conditions}
\label{sec:all-mixed-covariances}

The supplied continuation \emph{Full-support reconstruction: the missing
cohomology, the actual prime degree, and complete mixed-sign calculations}
calculates a two-factor example over $\Z/3$ and $\Z/5$. We retain its
actual amplitude maps and counting inner products and calculate the
entire spectrum for every finite number of factors. This is a finite
receiver of the split-support maps. The formulas below do not identify
it with an adelic completion or replace the original theta space by it.

Let $R_i$ be finite nonzero rings, $N_i=|R_i|\ge2$, and let $L_i$ be
finite nontrivial bounded distributive support lattices, $l_i=|L_i|$,
for $1\le i\le r$. All formal basis vectors in $\C[G_{L_i}(R_i)]$
and $H_i=\C[R_i]$ are orthonormal. Write $e_i=\delta_0$ and
$c_i=N_i^{-1/2}\sum_{x\in R_i}\delta_x$. The original amplitude map is
\[
 \pi_i\delta_{(a,1)}=\delta_a,\qquad
 \pi_i\delta_{z_\lambda}=\delta_0,\qquad
 \pi_i\pi_i^*=I+(l_i-1)e_ie_i^*.
 \tag{MG1}\label{mg:amplitude}
\]
The overlapping notation $(0,1)=z_1$ denotes one basis element.
There are exactly $l_i$ preimages of the amplitude zero and one
preimage of every nonzero amplitude, proving the last identity.

Put $H=\bigotimes_iH_i$, $e=\bigotimes_ie_i$, $c=\bigotimes_ic_i$,
$N=\prod_iN_i$, and $V=\{e,c\}^\perp$. Let $P$ be the orthogonal
projection onto $V$, $\pi=\bigotimes_i\pi_i$, and $T=P\pi$ with
codomain $V$. Its exact covariance is
\[
 C=TT^*=PDP|_V,\qquad
 D=\bigotimes_i\bigl(I+(l_i-1)e_ie_i^*\bigr).
 \tag{MG2}\label{mg:covariance}
\]
Both global constraints are present once. The operator $D$ is strictly
positive; hence $C$ is invertible on $V$ and $T$ is onto. The unique
minimum-norm lift of $v\in V$ is $T^*C^{-1}v$, because it has image
$v$ and is orthogonal to $\ker T$. Its squared norm is $v^*C^{-1}v$.
Thus $C^{-1}$ is exactly the attained quotient metric.

For each proper subset $S\subsetneq\{1,\ldots,r\}$ define
\[
 H_S=\bigotimes_{i\in S}\C e_i\ \otimes
            \bigotimes_{i\notin S}e_i^\perp,
 \quad q_S=\prod_{i\notin S}(N_i-1),
 \quad \alpha_S=\prod_{i\in S}l_i.
 \tag{MG3}\label{mg:blocks}
\]
The empty product is one. These are all the original mixed blocks
after removing the single all-zero vector $e$:
$e^\perp=\bigoplus_{S\subsetneq[r]}H_S$, with
$D|_{H_S}=\alpha_SI$ and $\dim H_S=q_S$.
Moreover $c_S=\operatorname{proj}_{H_S}c$ has
$\|c_S\|^2=q_S/N$. This follows by using
$\langle e_i,c_i\rangle=N_i^{-1/2}$ and multiplying all factor
norms. In particular every $c_S$ is nonzero and
$\sum_Sq_S=N-1$.

\begin{theorem}[Complete mixed covariance spectrum and determinant]
The characteristic polynomial on the original $(N-2)$-dimensional
space $V$ is exactly
\[
 \det(xI_V-C)=
 \prod_{S\subsetneq[r]}(x-\alpha_S)^{q_S-1}
 \frac{1}{N-1}
 \sum_{S\subsetneq[r]}q_S
       \prod_{\substack{T\subsetneq[r]\\T\ne S}}(x-\alpha_T).
 \tag{MG4}\label{mg:polynomial}
\]
This formula remains valid when different $\alpha_S$ coincide.
Its exact determinant and attained quotient determinant are
\[
 \det C=
 \left(\prod_{S\subsetneq[r]}\alpha_S^{q_S}\right)
 \frac{\sum_{S\subsetneq[r]}q_S/\alpha_S}{N-1},
 \qquad \det G_{\rm quotient}=(\det C)^{-1}.
 \tag{MG5}\label{mg:determinant}
\]
\end{theorem}
\begin{proof}
Inside $H_S$, the $(q_S-1)$-dimensional orthogonal complement of
$c_S$ belongs to $V$ and retains eigenvalue $\alpha_S$. What remains
is the compression of the diagonal matrix
$D_0=\operatorname{diag}(\alpha_S)$ to the hyperplane orthogonal
to $d=(\sqrt{q_S})_S$. In an orthonormal basis whose final vector
is $d/\sqrt{N-1}$, the determinant of the compressed block of
$xI-D_0$ is the final principal cofactor. For $x$ outside its
finite diagonal spectrum, the cofactor identity gives
\[
 \det(xI-D_0)\frac{d^*(xI-D_0)^{-1}d}{N-1}
 =\frac1{N-1}\sum_Sq_S\prod_{T\ne S}(x-\alpha_T).
\]
The identity holds for all $x$ because both expressions after
cancellation are polynomials. Multiplying the already separated
block factors proves MG4, including coincident eigenvalues.
Evaluating the same cofactor identity at $D_0$ itself proves
MG5, since every $\alpha_S>0$. The empty hyperplane when $r=1$
contributes determinant one, consistently with both formulas.
\end{proof}

The remaining eigenvalues can be located without cancelling a
mixed block. Between consecutive distinct values $\alpha_S$, the
secular function $\sum_Sq_S/(x-\alpha_S)$ has strictly negative
derivative, approaches $+\infty$ at the left endpoint from the
right, and approaches $-\infty$ at the right endpoint from the
left. It has exactly one root there. If a diagonal value occurs
at $h$ different subsets, its contribution in the compressed
remaining block is $h-1$. This follows either from MG4 or from
the $d$-orthogonal subspace of that equal-eigenvalue block.
All eigenvalues are positive and every multiplicity is specified.

For two factors, write $m=N_1$, $n=N_2$, $A=l_1$, $B=l_2$.
The extra two-dimensional characteristic polynomial is
\[
 \frac{(n-1)(x-B)(x-1)+(m-1)(x-A)(x-1)
 +(m-1)(n-1)(x-A)(x-B)}{mn-1}.
 \tag{MG6}\label{mg:two-factor}
\]
The other eigenvalues are $A$ with multiplicity $n-2$, $B$ with
multiplicity $m-2$, and $1$ with multiplicity
$(m-1)(n-1)-1$. Coincident values are combined using MG6.
In the supplied example $(m,n,A,B)=(3,5,2,2)$, MG6 equals
$(x-2)(x-11/7)$. Hence
\[
 \operatorname{spec}C=\{2^{[5]},(11/7)^{[1]},1^{[7]}\},
 \quad \det C=352/7,\quad\det G_{\rm quotient}=7/352.
 \tag{MG7}\label{mg:example}
\]
This recovers the supplied result from the complete factor-wise
amplitude map and determines all its higher-factor extensions.

There is an exact trace comparison for local versus global moment
conditions as well. For units $u_i\in R_i^\times$, let $U_i$
permute the basis by $x\mapsto u_ix$, and put
$f_i=|\{x\in R_i:(u_i-1)x=0\}|$. Both $e_i,c_i$ are fixed.
The tensor product action preserves $V$ and
$V_{\rm loc}=\bigotimes_i\{e_i,c_i\}^\perp\subseteq V$.
Their orthogonal complement inside $V$ is also invariant, and
\[
 \begin{aligned}
 \operatorname{Tr}_V U&=\prod_i f_i-2,\\
 \operatorname{Tr}_{V_{\rm loc}}U&=\prod_i(f_i-2),\\
 \operatorname{Tr}_{V\ominus V_{\rm loc}}U
 &=\prod_i f_i-2-\prod_i(f_i-2),\\
 \dim(V\ominus V_{\rm loc})&=N-2-\prod_i(N_i-2).
 \end{aligned}
 \tag{MG8}\label{mg:trace}
\]
The trace of each permutation is its fixed-point count; subtracting
its two fixed vectors gives the local formula. Tensoring and then
subtracting the global two-dimensional fixed space proves all four
identities. Multiplication by $2$ on $\Z/3$ and $\Z/5$ has
$f_1=f_2=1$, giving traces $-1,+1,-2$ on the three respective
spaces and dimension ten for the last. Because each $U_i$ fixes
$e_i$, it commutes with $D$, $P$, $C$ and the exact minimum-norm
section. Thus the trace comparison and the metric calculation
refer to the same original receiver.

\begin{figure}[ht]
\centering
\begin{tikzcd}[column sep=large,row sep=large]
 \bigotimes_i\C[G_{L_i}(R_i)]\arrow[r,"\pi"]
 &H\arrow[r,"P"]&V\\
 &&V\arrow[ull,bend left=15,"T^*C^{-1}"']\arrow[u,equal]
\end{tikzcd}
\caption{The original finite amplitude map retains every support
label before the two global moment conditions are imposed.
Equations MG1--5 calculate its complete covariance and the attained
minimum-norm section. Proper mixed blocks in MG3 survive; replacing
them by separate local two-condition spaces removes the trace and
dimension computed in MG8. The supplied two-factor example is
MG7; the all-factor formulas are derived here.}
\end{figure}
